% IEEE conference template — HPML Spring 2026 Final Report
% Team KV: Mamidibathula · Morisetty · Wang · Tian
% Columbia University
\documentclass[conference]{IEEEtran}

\usepackage{amsmath,amssymb,amsfonts}
\usepackage{textcomp}
\usepackage{xcolor}
\usepackage{booktabs}
\usepackage{multirow}
\usepackage{graphicx}
\usepackage{hyperref}
\usepackage{url}
\usepackage{microtype}

\hypersetup{colorlinks=true, linkcolor=blue, citecolor=blue, urlcolor=blue}

% ----------------------------------------------------------------
\begin{document}

\title{Comparative Evaluation of KV Cache Optimization Strategies\\
for Efficient LLM Inference: Quantization, Sparse Attention,\\
Eviction, and Latent Projection}

\author{%
\IEEEauthorblockN{%
  Rishika Mamidibathula\textsuperscript{\dag},
  Suhas Morisetty\textsuperscript{*},
  Ziheng Wang\textsuperscript{*},
  Tony Tian\textsuperscript{*}}
\IEEEauthorblockA{%
  \textsuperscript{*}Department of Computer Science,
  \textsuperscript{\dag}Department of Data Science,
  Columbia University\\
  COMS E6998: High Performance Machine Learning, Spring 2026\\
  \{rm4318, vm2825, zw3153, jt3640\}@columbia.edu}
}

\maketitle

% ----------------------------------------------------------------
\begin{abstract}
Autoregressive large language model (LLM) decoding is bounded by
key-value (KV) cache memory bandwidth: at batch size 32 and 4K context,
Llama-2-7B requires $\approx$65\,GB for KV alone, nearly saturating an
H100 80\,GB GPU.  We present a controlled, apples-to-apples comparison
of five KV-cache optimization techniques spanning four compression
families---quantization (KIVI 2/4-bit), sparse dynamic selection (TopK /
TokenSelect), one-shot eviction (SnapKV), and latent projection
(TransMLA)---plus a quantized-sparse hybrid (KIVI$\times$TopK).  All
methods share the same model weights, tokenizer, greedy decode loop,
prompts, seeds, and GPU (NVIDIA H100 80\,GB HBM3) so that differences
are attributable to the KV policy only.  Key findings: KIVI 4-bit is a
practical free lunch---lossless LongBench quality (0.292 vs.\ 0.291),
2.3$\times$ smaller KV cache, and 1.93$\times$ decode throughput at
BS\,=\,32 with the maximum single-GPU throughput improving 3.1$\times$
(1{,}517 vs.\ 497\,tok/s).  SnapKV marginally \emph{improves} overall
quality (+0.3\%) with zero per-step overhead.  TopK K\,=\,1024
matches baseline quality while achieving +20\% on TriviaQA via
noise suppression.  MLA delivers a consistent 3.9$\times$ KV reduction
and 1.3$\times$ decode speedup even at batch size 1, but the available
checkpoint exhibits a large quality regression requiring retraining.
No single method dominates all three axes simultaneously; the optimal
choice depends on the deployment's batch size, context length, and
quality tolerance.
\end{abstract}

\begin{IEEEkeywords}
KV cache, LLM inference, quantization, sparse attention, token eviction,
multi-head latent attention, throughput, memory efficiency
\end{IEEEkeywords}

% ================================================================
\section{Introduction}
% ================================================================

Large language model inference operates in two phases: prefill, which
processes the entire input prompt in a single forward pass, and decode,
which generates one token per step.  During decode, all past key and
value tensors must be read from GPU HBM at every step, making decode
strongly memory-bandwidth bound.  For Llama-2-7B \cite{llama2} with 32
transformer layers, 32 attention heads of dimension 128, FP16 precision,
4096-token context, and batch size 32, the KV cache alone occupies
$2 \times 32 \times 32 \times 128 \times 4096 \times 32 \times 2 \approx 65$\,GB
---nearly the full capacity of an H100 80\,GB GPU \cite{survey}.

Research on KV cache reduction has produced four principal compression
families:
(1)~\emph{Quantization}---reduce bits per element while keeping all
tokens \cite{kivi,kvquant};
(2)~\emph{Sparse selection}---attend to a dynamically chosen subset of
tokens per step \cite{tokenselect,rocketkv};
(3)~\emph{Eviction}---permanently discard low-importance tokens after
prefill \cite{h2o,snapkv,adakv};
(4)~\emph{Latent projection}---store a compressed low-rank
representation instead of full K/V tensors \cite{deepseekv2,transmla}.

Published evaluations rarely share identical experimental conditions,
making cross-family comparisons unreliable.  This work addresses that gap
with a unified benchmark harness that holds model, hardware, decode loop,
prompts, and seeds constant across all methods.

Our main contributions are:
(i) a single-codebase benchmark of five KV optimization techniques
including a quantized-sparse hybrid;
(ii) kernel-level analysis explaining crossover behavior (KIVI and TopK
are \emph{slower} than baseline at low batch sizes but faster at high
batch sizes or long context);
(iii) empirical evidence that KIVI 4-bit and SnapKV 0.4 provide
lossless or better LongBench quality;
(iv) full reproducibility via Modal cloud scripts and a public W\&B
dashboard at
\url{https://wandb.ai/rm4318-columbia-university/kv-cache-hpml}.

% ================================================================
\section{Literature Review}
% ================================================================

\textbf{KV Quantization.}
KIVI \cite{kivi} stores keys with per-channel scales (capturing the
dominant channel-direction outliers) and values with per-token grouped
quantization (group size 32), retaining the last 128 tokens in FP16 to
preserve short-range attention fidelity.  At 2-bit it achieves roughly
$8\times$ cache compression; at 4-bit, $4\times$.  KVQuant
\cite{kvquant} extends this with per-channel smoothing and non-uniform
codebooks, reaching near-lossless 1-bit compression at the cost of
greater overhead.

\textbf{Sparse Token Selection.}
TokenSelect \cite{tokenselect} (the basis of our TopK implementation)
maintains the full FP16 cache but at each decode step selects only the
top-$K$ tokens for attention, fusing the scoring kernel in Triton.  A
cosine-similarity reuse cache avoids re-scoring when consecutive queries
are similar.  RocketKV \cite{rocketkv} adds a two-stage compression
pipeline; RefreshKV \cite{refreshkv} periodically refreshes evicted
tokens.

\textbf{Eviction-Based Methods.}
H2O \cite{h2o} identifies ``heavy hitter'' tokens by cumulative
attention mass and evicts the rest.  SnapKV \cite{snapkv} performs a
single eviction after prefill by voting over a small observation window
of recent queries, producing a fixed-size compressed cache with no
per-step overhead.  Ada-KV \cite{adakv} allocates the budget
adaptively per head.

\textbf{Latent KV / Architectural Compression.}
DeepSeek-V2 \cite{deepseekv2} introduced Multi-head Latent Attention
(MLA), which stores a compressed latent vector $c_t$ per token and
expands to full K/V during attention via a learned up-projection, using
Attention Absorption to amortize the projection cost.  TransMLA
\cite{transmla} retrofits this structure onto an existing Llama-2
checkpoint.

% ================================================================
\section{Methodology}
% ================================================================

\subsection{Model and Hardware}

All experiments use \texttt{meta-llama/Llama-2-7b-chat-hf}
(32 layers, 32 heads, head dim 128, FP16 weights) on a single NVIDIA
H100 80\,GB HBM3 GPU provided by Modal cloud compute.  Software stack:
PyTorch 2.4.1, HuggingFace Transformers 4.44.2, CUDA 12.1, Triton 3.0,
flash-attn 2.5.  The MLA method additionally uses
\texttt{botsxc/llama2-7b-chat-mla-2048-8} \cite{transmla} (freqfold-8,
2048-dim latent, calibrated at 256-token context, no finetuning).

\subsection{Evaluation Datasets and Metrics}

\textbf{LongBench quality} \cite{longbench}: 6 long-context tasks
$\times$ 20 examples = 120 examples per method.  Tasks: \texttt{qasper}
(scientific QA, F1), \texttt{multifieldqa\_en} (multi-document QA, F1),
\texttt{triviaqa} (open-domain QA, F1), \texttt{2wikimqa} (multi-hop QA,
F1), \texttt{multi\_news} (summarization, ROUGE-L), \texttt{lcc} (code
completion, edit similarity).  Input truncated to 4096 tokens using
middle truncation (KIVI paper convention).

\textbf{Throughput benchmarks}: synthetic fixed-size batches; prefill
length 1024 tokens, generation 512 tokens, batch sizes 1--256.  Five
batches per config; first excluded as warmup; mean of four reported.

\textbf{Long-context benchmarks}: prefill length 32{,}768 tokens,
generation 512 tokens, batch size 1, using TopK-Flash and the
KIVI$\times$TopK hybrid.

\subsection{Profiling Tools}

\textbf{Memory:} \texttt{torch.cuda.max\_memory\_allocated()} reset
before each batch; measures peak, not steady-state.
\textbf{Timing:} \texttt{torch.cuda.synchronize()} + \texttt{time.perf\_counter()} for TTFT (time to first token, i.e., prefill) and decode throughput (tokens per second, steps 2--$G$).
\textbf{Experiment tracking:} Weights \& Biases; all runs at
\url{https://wandb.ai/rm4318-columbia-university/kv-cache-hpml}.
\textbf{Kernel profiling:} \texttt{torch.profiler} with Chrome trace
export for operator-level analysis.

\subsection{Method Implementations}

All methods implement a shared \texttt{MethodWrapper} interface in
\texttt{methods/} and are invoked through an identical generate loop in
\texttt{benchmark/runner.py}.

\textbf{Baseline (FP16):} Standard dense attention with no KV
modification; serves as the reference.

\textbf{KIVI \cite{kivi}:}  Keys quantized per-channel to 2 or 4 bits
(grouping across the sequence dimension per head); values quantized
per-token with group size 32.  The last \texttt{residual\_length}=128
tokens are retained in FP16.  Quantization uses a custom Triton kernel
(\texttt{triton\_quantize\_and\_pack}) that fuses INT4 packing; GEMV
dequantizes on-the-fly via \texttt{cuda\_bmm\_fA\_qB\_outer}.  The KV
cache is stored as a 9-tuple (quantized blocks + scales + zeros per K
and V + FP16 residual buffer) per layer.

\textbf{TopK / TokenSelect \cite{tokenselect}:} The full FP16 cache is
retained.  Three-region attention budget: sink window (first 128
tokens), dynamically-selected top-$K$ middle tokens, and local window
(last 512 tokens).  Per-decode-step scoring fuses $\mathbf{Q}\mathbf{K}^\top$
across all 32 heads via Triton kernel \texttt{\_topk\_score\_kernel}
$\rightarrow$ softmax $\rightarrow$ per-head sum $\rightarrow$
\texttt{torch.topk}.  A cosine similarity cache (threshold 0.9) skips
re-scoring when consecutive queries are similar.  RoPE correction passes
\texttt{true\_seq\_length} to each decode step.

\textbf{TopK-Flash:} Variant of TopK using a paged KV pool (256-token
pages) + \texttt{flash\_attn\_with\_kvcache} for zero-copy selective
attention.  Avoids the 52\,GB contiguous allocation required by standard
flash-attn at 32K context.

\textbf{SnapKV \cite{snapkv}:} After prefill, averages attention scores
over the last 32 query positions per head (the observation window),
applies max-pool (kernel size 7) to smooth noise, then retains
$\lfloor S \times 0.4 \rfloor$ top-voted tokens per head permanently.
Decode runs over the fixed reduced cache with zero per-step overhead.
Budget ratio 0.4; observation window 32.

\textbf{MLA (TransMLA \cite{transmla}):} Stores a compressed latent
vector $c_t \in \mathbb{R}^{d_c}$ per token (here $d_c = 2048$,
$4\times$ compression vs.\ $d_k \cdot n_h = 4096$).  At decode time,
$c_t$ is up-projected to full K/V by learned matrices; Attention
Absorption folds the up-projection into QKV attention.

\textbf{KIVI$\times$TopK Hybrid:} Combines KIVI (INT4 packed storage)
with sparse attention: prefill uses \texttt{F.scaled\_dot\_product\_attention}
(flash-attn backed, $O(T)$ memory) to avoid 32K OOM; decode scores all
tokens via \texttt{cuda\_bmm\_fA\_qB\_outer} over INT4-packed K, selects
$n_\text{sink}+K+n_\text{local}$ positions, then gathers packed INT4 V
rows directly for the attention GEMV without an intermediate FP16 write.

% ================================================================
\section{Experimental Results}
% ================================================================

\subsection{LongBench Quality}

Table~\ref{tab:quality} reports per-task and overall LongBench scores
for all methods.  KIVI 4-bit and TopK K=1024 both achieve 0.292
(+0.3\% and +0.4\% over baseline 0.291) --- within evaluation noise for
20 examples per task but consistent across runs.  SnapKV 0.4 reaches
0.295 (+0.3\%).

\begin{table}[!t]
\centering
\caption{LongBench quality (F1 / ROUGE-L / edit-sim), 6 tasks $\times$ 20 ex.
Tasks: (1) qasper, (2) mfqa, (3) triviaqa, (4) 2wikimqa,
(5) mnews, (6) lcc.}
\label{tab:quality}
\setlength{\tabcolsep}{3.5pt}
\begin{tabular}{@{}lrrrrrrr@{}}
\toprule
\textbf{Method} & \textbf{(1)} & \textbf{(2)} & \textbf{(3)} & \textbf{(4)} & \textbf{(5)} & \textbf{(6)} & \textbf{Avg} \\
\midrule
Baseline (FP16)       & .172 & .380 & .644 & .083 & .181 & .287 & .291 \\
KIVI 4-bit            & .194 & .360 & .643 & .083 & .188 & .286 & \textbf{.292} \\
KIVI 2-bit            & .133 & .382 & .608 & .050 & .191 & .277 & .273 \\
TopK K=2048           & .145 & .329 & .731 & .075 & .178 & .282 & .290 \\
TopK K=1024           & .153 & .281 & .773 & .089 & .183 & .273 & \textbf{.292} \\
TopK K=512            & .145 & .237 & .751 & .050 & .178 & .281 & .274 \\
SnapKV 0.4            & .153 & .387 & .687 & .083 & .181 & .277 & \textbf{.295} \\
MLA (B=1)*            & .036 & .101 & .064 & .043 & .103 & .140 & .081 \\
\bottomrule
\end{tabular}
\vspace{1pt}
\footnotesize{* MLA baseline truncated to 2048 tokens vs.\ 4096 for others.}
\end{table}

Per-task patterns reveal complementary failure modes:
\emph{KIVI 2-bit and TopK K=512 both converge to overall score
$\approx$0.274 via opposite mechanisms.}  KIVI's quantization noise
most hurts qasper (−22\%) and 2wikimqa (−40\%), tasks requiring precise
recall of sparse mid-context facts.  TopK's sparse selection most hurts
multifieldqa\_en (−38\% at K=512) because answer spans at depth $> n_\text{sink} + n_\text{local}$ are pruned.  Conversely, \emph{TopK
beats baseline on TriviaQA across all K} (+13\% to +20\%): the answer
appears in multiple phrasings; sparse selection drops distracting copies
and focuses attention.  SnapKV's one-shot vote actively aligns selection
with answer-bearing positions, boosting triviaqa (+4\%) and
multifieldqa (+0.7\%) while only hurting qasper (−2\%) where
evidence spans fall outside the 32-token observation window.

\subsection{Throughput and Memory}

Table~\ref{tab:throughput} shows the KIVI throughput sweep at prefill
1024, gen 512.  Key observations:

\begin{table}[!t]
\centering
\caption{Decode throughput (tok/s) and peak GPU memory vs.\ batch size.
Prefill 1024, gen 512. ``—'' = OOM.}
\label{tab:throughput}
\setlength{\tabcolsep}{4pt}
\begin{tabular}{@{}rrrrrrr@{}}
\toprule
 & \multicolumn{3}{c}{\textbf{Decode tok/s}} & \multicolumn{3}{c}{\textbf{Peak mem (GB)}} \\
\cmidrule(lr){2-4} \cmidrule(lr){5-7}
\textbf{BS} & \textbf{Base} & \textbf{KIVI-2b} & \textbf{Spdup} & \textbf{Base} & \textbf{KIVI-2b} & \textbf{MLA} \\
\midrule
1   &  70 &   30 & 0.43$\times$ & 15.1 & 14.0 & 13.8 \\
2   & 135 &   61 & 0.45$\times$ & 16.7 & 14.4 &  —   \\
4   & 242 &  123 & 0.51$\times$ & 20.0 & 15.3 &  —   \\
8   & 351 &  243 & 0.69$\times$ & 26.4 & 17.1 &  —   \\
16  & 438 &  480 & 1.10$\times$ & 39.3 & 20.7 &  —   \\
32  & 497 &  957 & 1.93$\times$ & 65.0 & 27.8 &  —   \\
64  &  —  & 1295 & $\infty$     &  —   & 42.1 &  —   \\
128 &  —  & 1517 & $\infty$     &  —   & 70.7 &  —   \\
\bottomrule
\end{tabular}
\end{table}

\textbf{Crossover at BS=16.}  KIVI 2-bit decode throughput is
$0.4$--$0.7\times$ baseline at BS 1--8 but crosses over at BS=16
(480 vs.\ 438\,tok/s, 1.10$\times$) and reaches 1.93$\times$ at
BS=32.  The crossover arises because at small batch the decode
$\mathbf{Q}\mathbf{K}^\top$ is a GEMV: INT4 GEMV does not map to
tensor core units and is $\approx$2.5$\times$ slower than FP16 GEMV.
At large batch the operation becomes a near-GEMM where HBM bandwidth
saving (the KV cache is $2.3\times$ smaller at BS=32: 27.8 vs.\ 65\,GB)
outweighs dequantization cost.

\textbf{Capacity extension.}  Baseline OOMs at BS=64; KIVI 2-bit
reaches BS=128 at 70.7\,GB and achieves 1{,}517\,tok/s—$3.1\times$
the maximum ever achievable with the baseline on this GPU.

\textbf{TTFT overhead.}  KIVI TTFT is 18--66\% worse across batch sizes
(e.g., 998 vs.\ 1659\,ms at BS=32) because the full KV cache is
quantized and packed at the end of prefill---a cost proportional to
$B \times S$.  For latency-sensitive workloads (interactive chat), this
is a real penalty; for throughput-oriented offline inference, KIVI
wins decisively.

\textbf{MLA at B=1.}  MLA achieves 43--44\,tok/s vs.\ 34\,tok/s
baseline (1.26--1.32$\times$) despite batch size 1.  This confirms
that even single-request decode is partially HBM-bandwidth bound; the
3.9$\times$ KV reduction reduces per-step HBM traffic proportionally.
TTFT is also lower with MLA at shorter prefill lengths (39 vs.\ 52\,ms
at $p$=512) due to the smaller working set.

\subsection{Long-Context Results (32K)}

\begin{table}[!t]
\centering
\caption{Long-context results. Prefill 32{,}768 tokens, gen 512, BS=1.}
\label{tab:longctx}
\begin{tabular}{@{}lrrrr@{}}
\toprule
\textbf{Method} & \textbf{TTFT (ms)} & \textbf{Spdup} & \textbf{Peak mem} & \textbf{Saving} \\
\midrule
Baseline (FP16)     & 2597 & 1.00$\times$ & 52.1\,GB & — \\
TopK-Flash          & 1761 & 1.47$\times$ & 42.1\,GB & −19\% \\
SnapKV 0.4          &  —   & —            & 26.0\,GB & −50\% \\
KIVI$\times$TopK 4b &  —   & 0.64$\times$ & 25.8\,GB & −51\% \\
\bottomrule
\end{tabular}
\vspace{1pt}
\footnotesize{Baseline OOMs at 64K; TopK-Flash succeeds at 65.8\,GB.}
\end{table}

\textbf{TopK-Flash.}  At 32K context the 256-token paged KV pool
eliminates the 52\,GB contiguous memory allocation that standard
flash-attn requires, reducing TTFT by 33\% and peak memory by 19\%.
The gain is not from sparsity per se but from allocation and
memory-access-pattern improvement.  Baseline OOMs at 64K; TopK-Flash
handles it at 65.8\,GB.

\textbf{KIVI$\times$TopK hybrid.}  At 32K, combining INT4 storage with
sparse attention cuts peak memory 51\% (52.1$\rightarrow$25.8\,GB) but
decode speed drops 36\% (21.4$\rightarrow$13.7\,tok/s).  The slowdown
arises because $\mathbf{Q}\mathbf{K}^\top$ must still score all 32K
tokens every step before top-$K$ can be applied---the scoring cost is
identical to the baseline, plus additional topk + sort + gather
overhead.  LongBench quality remains at $\approx$0.297 (within noise of
0.291 baseline).

% ================================================================
\section{Discussion}
% ================================================================

\subsection{Deployment Recommendations}

Table~\ref{tab:tradeoff} summarizes the operating point of each method.

\begin{table}[!t]
\centering
\caption{Deployment fit summary.}
\label{tab:tradeoff}
\setlength{\tabcolsep}{3pt}
\begin{tabular}{@{}lllll@{}}
\toprule
\textbf{Method} & \textbf{KV ratio} & \textbf{Quality} & \textbf{Best regime} \\
\midrule
KIVI 4-bit  & 4$\times$  & Lossless (.292)          & Batch serving $\geq$ BS=16 \\
KIVI 2-bit  & 8$\times$  & −6.2\% (.273)            & Max-batch, accept noise \\
TopK K=1024 & 1$\times$† & Lossless (.292)          & Single-req, long context \\
SnapKV 0.4  & 2.5$\times$& +0.3\% (.295)            & Any; quality-neutral \\
MLA (B=1)   & 3.9$\times$& −65\%* needs retrain     & After finetuning at $\geq$4K \\
\bottomrule
\end{tabular}
\vspace{1pt}
\footnotesize{† TopK does not reduce KV storage; * truncation mismatch may inflate gap.}
\end{table}

\textbf{KIVI 4-bit} is the headline recommendation for high-throughput
batch inference: it costs nothing in quality, extends the maximum
serviceable batch size from 32 to 128, and multiplies peak throughput
3.1$\times$ on a single H100.  The only cost is TTFT latency
(+18--66\%), which matters for interactive applications.

\textbf{SnapKV 0.4} is the best choice when quality must be preserved
and per-step overhead must be zero.  Its 2.5$\times$ KV reduction comes
``free'' during serving because the compressed cache is built once
after prefill.

\textbf{TopK K=1024} is suited to single-request, long-context
scenarios ($>$32K) where the paged variant (TopK-Flash) provides real
TTFT and memory gains.  At short context and low batch, scoring overhead
dominates.

\textbf{MLA} shows the theoretical ceiling of latent projection: 3.9$\times$
KV reduction and 1.3$\times$ decode speedup are repeatable benefits
even at B=1, but realizing this on a production-quality checkpoint
requires finetuning with a longer context window and higher latent rank.

\subsection{Kernel-Level Observations}

\textbf{INT4 GEMV inefficiency.}  \texttt{cuda\_bmm\_fA\_qB\_outer}
performs per-GEMV INT4$\rightarrow$FP16 dequantization.  H100 tensor
cores target INT8/FP16 GEMM; INT4 GEMV falls back to scalar or
vector-INT4 paths and runs $\approx$2.5$\times$ slower than native FP16
GEMV at BS=1.  The bandwidth saving only exceeds this overhead when
decode approaches the GEMM regime at BS$\geq$16.

\textbf{KIVI residual buffer hot path.}  At short generation lengths
($< \texttt{residual\_length} \times n_\text{layers}$), nearly all
attention concentrates in the FP16 residual; KIVI provides minimal
storage benefit.  This is also why KIVI TTFT is $3\times$ slower:
full-prefill quantization is paid unconditionally.

\textbf{Triton fused scoring.}  \texttt{\_topk\_score\_kernel} fuses
$\mathbf{Q}\mathbf{K}^\top$ + softmax + head-sum into a single kernel
launch, reducing Python-dispatch overhead 2--3$\times$ vs.\ sequential
PyTorch ops.  At 1K--4K sequence lengths the critical path remains HBM
bandwidth, not compute.

\subsection{Limitations}

\emph{Single model architecture.}  All results are specific to
Llama-2-7B with multi-head attention (32 heads, 32 KV heads).  Models
with grouped-query attention (Llama-3, Mistral) have fewer KV heads,
changing the compression profile.

\emph{Small evaluation set.}  LongBench results are based on 20
examples per task.  Some per-task deltas (e.g., qasper KIVI 2-bit
$-22\%$) may partially reflect small-$n$ variance rather than a
systematic effect, although the direction is consistent with the
expected behavior.

\emph{MLA checkpoint mismatch.}  The available TransMLA checkpoint was
calibrated at 256-token context without task finetuning.  A properly
trained MLA checkpoint would likely recover much of the 65\% quality gap.

\emph{Hybrid scoring cost.}  The KIVI$\times$TopK hybrid's memory win
(51\%) comes with a decode penalty (−36\%) because scoring must precede
selection.  Future work using a hierarchical or lazy-scoring approach
could break this dependency.

% ================================================================
\section{Conclusion}
% ================================================================

We conducted a controlled benchmark of five KV cache optimization
methods across four compression families on Llama-2-7B.  KIVI 4-bit
provides lossless quality with 4$\times$ KV compression and 1.93$\times$
throughput at BS=32, making it the strongest practical choice for
batch-oriented deployment today.  SnapKV marginally improves LongBench
overall score with zero runtime overhead.  TopK K=1024 matches baseline
quality while gaining on noise-heavy retrieval tasks.  MLA's per-token
latent representation yields consistent 3.9$\times$ KV reduction and
1.3$\times$ decode speedup at B=1, but requires a properly trained
checkpoint for production quality.  No single method dominates all
axes; practitioners should select based on whether batch size,
context length, or output fidelity is the primary constraint.

Future work should address: (i)~batched decode for TopK to amortize the
per-step scoring cost and achieve a KIVI-style crossover at BS$\geq$8;
(ii)~MLA retraining at 4K--8K context with rank 16--32 to preserve the
3.9$\times$ compression gain while recovering QA quality;
(iii)~lazy or hierarchical scoring in the KIVI$\times$TopK hybrid to
decouple the memory win from the decode penalty; and
(iv)~cross-method composition (e.g., SnapKV eviction + KIVI
quantization on the retained tokens) to stack savings from orthogonal
axes.

\appendix
\section*{AI Use Disclosure}

Per the HPML AI Use Policy.  We used Claude (Anthropic) and GitHub
Copilot for: debugging CUDA OOM errors, clarifying Triton kernel
semantics, and polishing report prose.  All reported experimental
numbers were produced by running our own code on Modal H100 containers.
All code was reviewed and validated by team members.

% ================================================================
\bibliographystyle{IEEEtran}

\begin{thebibliography}{14}

\bibitem{survey}
J.~Jiang, P.~Yang, R.~Zhang, and F.~Liu,
``Towards efficient large language model serving: A survey on
system-aware KV cache optimization,''
\textit{TechRxiv}, 2026.

\bibitem{llama2}
H.~Touvron \textit{et al.},
``Llama 2: Open foundation and fine-tuned chat models,''
\textit{arXiv:2307.09288}, 2023.

\bibitem{kivi}
Z.~Liu \textit{et al.},
``KIVI: A tuning-free asymmetric 2-bit quantization for KV cache,''
in \textit{Proc.\ ICML}, 2024.

\bibitem{kvquant}
C.~Hooper \textit{et al.},
``KVQuant: Towards 10 million context length LLM inference with KV
cache quantization,''
in \textit{NeurIPS}, 2024.

\bibitem{tokenselect}
W.~Wu \textit{et al.},
``TokenSelect: Efficient long-context inference via dynamic token-level
KV cache selection,''
in \textit{Proc.\ EMNLP}, 2025.

\bibitem{rocketkv}
P.~Behnam \textit{et al.},
``RocketKV: Accelerating long-context LLM inference via two-stage KV
cache compression,''
in \textit{Proc.\ ICML}, 2025.

\bibitem{refreshkv}
F.~Xu, T.~Goyal, and E.~Choi,
``RefreshKV: Updating small KV cache during long-form generation,''
in \textit{Proc.\ ACL}, 2025.

\bibitem{h2o}
Z.~Zhang \textit{et al.},
``H2O: Heavy-hitter oracle for efficient generative inference of large
language models,''
in \textit{NeurIPS}, 2023.

\bibitem{snapkv}
Y.~Li \textit{et al.},
``SnapKV: LLM knows what you are looking for before generation,''
in \textit{NeurIPS}, 2024.

\bibitem{adakv}
Y.~Feng \textit{et al.},
``Ada-KV: Optimizing KV cache eviction by adaptive budget allocation,''
in \textit{NeurIPS}, 2025.

\bibitem{deepseekv2}
DeepSeek-AI,
``DeepSeek-V2: A strong, economical, and efficient mixture-of-experts
language model,''
\textit{arXiv:2405.04434}, 2024.

\bibitem{transmla}
Y.~He \textit{et al.},
``TransMLA: Multi-head latent attention is all you need,''
\textit{arXiv:2502.07864}, 2025.

\bibitem{xkv}
C.-C.~Chang \textit{et al.},
``xKV: Cross-layer SVD for KV-cache compression,''
\textit{arXiv:2503.18893}, 2025.

\bibitem{longbench}
Y.~Bai \textit{et al.},
``LongBench: A bilingual, multitask benchmark for long context
understanding,''
\textit{arXiv:2308.14508}, 2023.

\end{thebibliography}

\end{document}
